% Generated draft: requires mathematical and source review.
\documentclass[11pt]{article}
\usepackage[margin=25mm]{geometry}
\usepackage{fontspec}
\setmainfont{Latin Modern Roman}
\usepackage{amsmath,amssymb,mathtools,unicode-math}
\setmathfont{Latin Modern Math}
\usepackage{enumitem,needspace,xurl,hyperref}
\setlength{\parindent}{0pt}
\setlength{\parskip}{5pt}
\setlength{\emergencystretch}{3em}
\newcommand{\N}{\mathbb{N}}
\newcommand{\Z}{\mathbb{Z}}
\newcommand{\Q}{\mathbb{Q}}
\newcommand{\R}{\mathbb{R}}
\newcommand{\C}{\mathbb{C}}
\newcommand{\F}{\mathbb{F}}
\newcommand{\A}{\mathbb{A}}
\newcommand{\PP}{\mathbb{P}}
\newcommand{\E}{\mathbb{E}}
\newcommand{\eps}{\varepsilon}
\newcommand{\dd}{\,\mathrm{d}}
\begin{document}
\section*{UnsolvedMath ID: 8800002}
\textbf{Source problem code:} catalogue entry
\subsection{Definitions and mathematical statement}
A bilinear pairing between abelian groups is a map $e:G_1\times G_2\to G_T$ that is a homomorphism in each argument. In additive source and multiplicative target notation,
\[
e(P+P',Q)=e(P,Q)e(P',Q),\qquad
e(P,Q+Q')=e(P,Q)e(P,Q').
\]
Consequently $e(aP,bQ)=e(P,Q)^{ab}$ for integers $a,b$. A common nondegeneracy condition requires that $e(P,Q)=1$ for every $Q$ imply $P=0$, and the corresponding condition in the other argument. These are explanatory standard pairing notions, not all requirements explicitly imposed by the source.

The problem is to find bilinear structures useful for cryptographic constructions. It supplies no specific groups, parameter families, algorithms, security assumption, efficiency bound, nondegeneracy convention, or cryptographic task. Finite cyclic prime-order groups and efficiently computable nondegenerate pairings are a common setting, but are not a restriction in the given prompt. Usefulness must be assessed for a specified construction and security model; bilinearity alone is not a cryptographic guarantee.

This is a research programme, not a claim that a particular pairing exists with universal security or that known pairing examples exhaust the search.
\subsection{Short English statement}
Identify bilinear structures that support useful cryptographic constructions, with appropriate algebraic, computational and security properties specified for each proposed application.
\subsection{Sources}
\begin{itemize}
\item Local UnsolvedMath registry, permanent ID 8800002, exact source code catalogue entry, ``Cryptographically useful bilinear structures''; attributed to Gerhard Frey, no proposed year supplied. Registry status partially solved, classification PARTIAL-PROGRESS.
\item \emph{IPAM Workshop on Number Theory and Cryptography---Open Problems} (2006), attachment 6219, PDF page 8, Problem 1, \url{https://www.ipam.ucla.edu/programs/workshops/workshop-i-number-theory-and-cryptography-open-problems/?tab=schedule}. Supplied extraction: pdf-and-legacy-ppt visual audit, accessed July 29, 2026. Its provenance says 19 surviving attachments, 17 PDFs and two legacy PPT decks, were inspected and repeated animations/solved definitions excluded; this was the source audit, not an inspection performed for this draft.
\item The supplied summary notes extensive Weil/Tate/ate pairing, divisor-class and Brauer-group developments, especially elliptic-curve genus-one work, but gives no exact references and says the broad request cannot be declared solved. No genus-one restriction is added.
\item Provenance flags current primary-literature status and redistribution terms NEEDS\_REVIEW. No PDF or PPT was accessed here.
\end{itemize}
\end{document}
